\documentclass[11pt,letterpaper]{article}
\usepackage[margin=1in,headheight=15pt]{geometry}
\usepackage[T1]{fontenc}
\usepackage{amsmath,amsthm}
\usepackage{newtxtext,newtxmath,microtype,mathtools}
\usepackage{booktabs,longtable,array,enumitem,fancyhdr}
\usepackage[dvipsnames]{xcolor}
\usepackage[most]{tcolorbox}
\usepackage{hyperref}
\hypersetup{colorlinks=true,linkcolor=MidnightBlue,citecolor=MidnightBlue,urlcolor=MidnightBlue,
 pdftitle={An exact gamma constant for A279619},pdfauthor={Research report prepared for Vladimir Reshetnikov}}
\usepackage{bookmark}
\setlength{\parskip}{4pt}\setlength{\parindent}{1.2em}
\setlength{\emergencystretch}{2em}\allowdisplaybreaks[2]
\pagestyle{fancy}\fancyhf{}
\fancyhead[L]{\small A279619: an exact gamma constant}
\fancyhead[R]{\small Research report}\fancyfoot[C]{\thepage}
\newtheorem{theorem}{Theorem}[section]
\newtheorem{proposition}[theorem]{Proposition}
\newtheorem{lemma}[theorem]{Lemma}
\newtheorem{corollary}[theorem]{Corollary}
\theoremstyle{remark}\newtheorem{remark}[theorem]{Remark}
\newcommand{\G}{\mathcal G_7}\newcommand{\Ast}{A_\star}
\newcommand{\Q}{\mathbb Q}\newcommand{\tht}{\vartheta}
\newcommand{\hyp}[4]{{}_2F_1\!\left(\begin{matrix}#1,#2\\#3\end{matrix};#4\right)}
\newtcolorbox{resultbox}{enhanced,breakable,colback=blue!3,colframe=MidnightBlue!65,
 boxrule=.5pt,arc=1pt,left=10pt,right=10pt,top=7pt,bottom=7pt}
\begin{document}
\begin{titlepage}
\vspace*{.4in}
{\large\sffamily\color{MidnightBlue} OEIS RESEARCH \quad / \quad LEVEL SEVEN}\par
\vspace{.3in}
{\huge\bfseries An exact gamma constant\par for an Ap\'ery-like sequence}\par
\vspace{.2in}
{\Large A279619, all-orders asymptotics,\par and its recessive companion}\par
\vspace{.3in}
{\large Research report prepared for Vladimir Reshetnikov}\par
{\large October 1, 2026}\par
\vspace{.3in}
\begin{resultbox}
\[
 a_n\sim\frac{\sqrt{3\pi}}{\Gamma(1/7)\Gamma(2/7)\Gamma(4/7)}
               \frac{27^n}{n^{3/2}}.
\]
The leading constant is evaluated exactly, the displayed OEIS asymptotic
is proved, and the expansion is extended to arbitrary fixed order with a
rigorous remainder. A rationally normalized companion exposes an
exponentially smaller alternating scale.
\end{resultbox}
\vspace{.2in}
\noindent\textbf{Abstract.}
We study the level-seven Ap\'ery-like sequence A279619, defined by a
second-order polynomial recurrence. Its OEIS entry displays a conjectural
asymptotic with a numerical constant. We prove that asymptotic and identify
the constant as a reciprocal product of three gamma values. A rational
parameter removes the square roots in the known hypergeometric generating
function, making its singular connection coefficient explicit. Classical
elliptic-period and Ramanujan identities reduce this coefficient to its
gamma-product form. Frobenius expansion and singularity analysis give all
algebraic correction terms, with rational coefficients and fixed-order
error bounds. We also construct a companion continued fraction, prove
exact rational bracketing, and determine the amplitude and every correction
term of its recessive solution. Applications include convolution powers,
A183204, eventual log-convexity, and Lambert-$W$ inversion with a separate
treatment of integer rounding. Exact symbolic certificates and rational
enclosures accompany numerical experiments. Classical inputs, previously
calculated terms, sequence-specific deductions, and unsettled questions
are explicitly distinguished.
\vfill
\noindent\small
Deliverables: TeX source, compiled PDF, executable verification program,
rational certificates, numerical tables, and a source audit.\\
These are conventional mathematical proofs with computer-algebra checks,
not Lean-certified theorems. Global priority is not claimed.
\end{titlepage}
\tableofcontents\newpage

\section{The target and the scope of the contribution}
The sequence begins
\[
1,\ 2,\ 22,\ 336,\ 6006,\ 117348,\ 2428272,\ 52303680,\ldots
\]
and satisfies
\begin{equation}\label{eq:rec}
(n+1)^2a_{n+1}=(26n^2+13n+2)a_n+3(3n-1)(3n-2)a_{n-1},\qquad n\ge1,
\end{equation}
with $a_0=1$, $a_1=2$. Throughout, set
\begin{equation}\label{eq:defs}
 A(z)=\sum_{n\ge0}a_nz^n,\qquad \rho=1/27,\qquad
 \G=\Gamma(1/7)\Gamma(2/7)\Gamma(4/7).
\end{equation}
This is A279619; $A(z)^2$ is the generating function of A183204
\cite{oeis279619,oeis183204}. Its connections with modular forms,
hypergeometric functions, and another Ap\'ery-like sequence make it a
structured target rather than an isolated numerical curiosity.

\subsection{What is already known}
The OEIS entry consulted on October 1, 2026 labels its leading asymptotic
as conjectural and supplies a numerical constant. O'Brien's thesis
\cite[\S8.4, pp.~63--67]{obrien} estimates that constant and formally
calculates the first three correction terms. Those three terms are
\emph{not new here}. The hypergeometric representation is also recorded
in OEIS in code attributed to V\'aclav Kote\v{s}ovec in 2018
\cite{oeis279619}. We prove and rationalize that representation before
using it analytically.

The contribution is not a new general method of asymptotic analysis.
It is an exact connection-constant evaluation, a rigorous all-orders
fixed-order expansion, and an explicit recessive companion with a proved
error amplitude. Standard machinery is used to close the displayed
asymptotic question.
\begin{center}
\begin{tabular}{p{.30\textwidth}p{.61\textwidth}}
\toprule
Item & Status in this report\\\midrule
Recurrence and hypergeometric function & Existing input; the latter
receives a rational differential-equation certificate.\\
Leading scale and first three corrections & Previously conjectured or
formally calculated; now given a remainder proof.\\
Gamma-product constant & Exact sequence-specific evaluation deduced
from classical special values.\\
Companion continued fraction & Constructed with exact bracketing and
all-orders recessive asymptotics; global priority not asserted.\\
Convolution and inverse expansions & Deductions with explicit formulas
and integer-rounding bounds.\\
Lucas congruences & Not resolved or used here.\\\bottomrule
\end{tabular}
\end{center}
A targeted search did not locate an earlier proof of this particular
constant formula. This is not an exhaustive priority search. In
particular, recent work of Cohen and Zudilin on continued fractions for
Chowla--Selberg quotients \cite{cohenzudilin} shows that the broader area
is active. We do not claim a first continued fraction for a gamma
quotient, an identification of the companion limit with such a quotient,
or an irrationality measure.

\subsection{Relation to ProveIt}
The repository's \texttt{Analysis/Transseries} project separates
asymptotic scales, flatness, power--logarithmic expressions, and inversion
from the Fabius-function development \cite{proveit}. This report follows
that separation: an analytic existence proof precedes exact coefficient
recursions and inverse expansions. No unverified repository theorem is
needed. The inspected snapshot was
\texttt{a5ad4ac2e2d4e642a6c015836a70492d40fcbc71}; inspection was not an
exhaustive review of archived reports.

\section{Principal results}
\begin{theorem}[Dominant asymptotics and the exact constant]\label{th:main}
Put
\begin{equation}\label{eq:C}
 C=\frac{\sqrt{3\pi}}{\G},\qquad \Ast=\frac{3\G}{8\pi^2}.
\end{equation}
Then $A(\rho)=\Ast$. For every fixed $M\ge0$,
\begin{equation}\label{eq:dominant}
 a_n=C\frac{27^n}{n^{3/2}}
 \left(\sum_{j=0}^M\frac{c_j}{n^j}+O_M(n^{-M-1})\right),
\end{equation}
where the $c_j$ are uniquely determined rational numbers with $c_0=1$.
In particular,
\begin{align}\label{eq:cvalues}
 c_1&=-\frac{215}{1008},&c_2&=-\frac{1265}{290304},&
 c_3&=\frac{4683055}{877879296},\notag\\
 c_4&=\frac{924908071}{3539609321472},&
 c_5&=-\frac{6088275239005}{3567926196043776}.
\end{align}
An exact triangular algorithm computes every coefficient.
\end{theorem}
The theorem also evaluates a convergent boundary sum:
\begin{equation}\label{eq:boundarysum}
 \sum_{n=0}^{\infty}\frac{a_n}{27^n}
 =\frac{3\Gamma(1/7)\Gamma(2/7)\Gamma(4/7)}{8\pi^2},
 \qquad C\Ast=\frac{3\sqrt3}{8\pi^{3/2}}.
\end{equation}
The gamma product in the boundary value cancels against its reciprocal
in the coefficient amplitude.

\begin{theorem}[An exact recessive companion]\label{th:companion}
Let $b_0=0$, $b_1=1$, and impose \eqref{eq:rec} on $b_n$ for $n\ge1$.
Then the limit
\begin{equation}\label{eq:L}
 L=\lim_{n\to\infty}\frac{b_n}{a_n}
 =0.46721176888437312394211005141309126172410854995971\ldots
\end{equation}
exists and satisfies
\begin{equation}\label{eq:bracket}
 \frac{b_{2m}}{a_{2m}}<L<\frac{b_{2m+1}}{a_{2m+1}},\qquad m\ge0.
\end{equation}
For $e_n=b_n-La_n$ and every fixed $M\ge0$,
\begin{equation}\label{eq:minimal}
 e_n=(-1)^{n+1}\frac{\sqrt3}{56\pi C}\,n^{-3/2}
 \left(\sum_{j=0}^M\frac{d_j}{n^j}+O_M(n^{-M-1})\right),
\end{equation}
where $d_0=1$ and
\begin{equation}\label{eq:dvalues}
 d_1=-\frac{57}{112},\quad d_2=\frac{655}{3584},\quad
 d_3=-\frac{13365}{401408},\quad d_4=-\frac{3627123}{179830784}.
\end{equation}
The $d_j$ are rational and computable to arbitrary order. Consequently,
\begin{equation}\label{eq:CFerror}
 L-\frac{b_n}{a_n}\sim(-1)^n\frac{\sqrt3}{56\pi C^2}\,27^{-n}.
\end{equation}
\end{theorem}
The two scales are $27^nn^{-3/2}$ and $(-1)^nn^{-3/2}$. The second is
exponentially small \emph{relative} to the first, not in absolute
magnitude. We do not assign an unproved canonical subdominant coefficient
to the asymptotic expansion of $a_n$ itself.

\section{From the recurrence to a controlled singularity}
\subsection{Positivity, growth, and the differential equation}
The recurrence immediately gives $a_n>0$. Induction also gives
$a_n\le27^n$, since
\[
26n^2+13n+2+\frac{3(3n-1)(3n-2)}{27}
=27n^2+12n+\frac{20}{9}<27(n+1)^2.
\]
Thus $A$ is analytic in $|z|<\rho$. Multiplying the recurrence by $z^n$
and summing, including $a_1=2a_0$, gives
\begin{equation}\label{eq:ODE}
 zD(z)A''+(1-39z-54z^2)A'-(2+6z)A=0,
 \qquad D(z)=(1-27z)(1+z).
\end{equation}
The coefficient of $z^n$ on the left is the recurrence residual, so the
analytic solution normalized by $A(0)=1$ is unique. Outside
$0,\rho,-1$, the equation is ordinary; analytic continuation along paths
avoiding these points follows from linear differential-equation theory.

\subsection{The local exponents}
In $s=1-27z$, the equation becomes
\begin{equation}\label{eq:sODE}
 s(1-s)(28-s)y''+(14-43s+2s^2)y'-\frac{20-2s}{9}y=0.
\end{equation}
The indicial equation is $14\alpha(2\alpha-1)=0$. Its exponents $0,1/2$
have nonintegral difference, hence Frobenius theory gives a convergent
local expansion
\begin{equation}\label{eq:puiseux}
 A(z)=U(s)+B\sqrt{s}\,V(s),\qquad V(0)=1,
\end{equation}
where $U,V$ are analytic at $0$. Branches are positive for $0<z<\rho$.
Below we show $B<0$, so the singularity is not removable.

The self-adjoint form provides a useful sign check:
\begin{equation}\label{eq:flux}
 \bigl(z\sqrt{D(z)}A'(z)\bigr)'
   =\frac{2+6z}{\sqrt{D(z)}}A(z),\qquad 0<z<\rho.
\end{equation}
Frobenius theory makes $A(\rho)$ finite. Positivity of the right side
implies that the limiting flux is positive, whereas \eqref{eq:puiseux}
identifies it as $-\tfrac12\sqrt{28/27}\,B$.

There is no other singularity on $|z|=\rho$. The chosen solution extends
to a dented disk of radius greater than $\rho$, cut to the right from
$\rho$. This supplies the global analytic condition for singularity
analysis, not just a formal local expansion \cite[Chapter VI]{flajolet}.

\section{A rational certificate for the generating function}
Set
\begin{equation}\label{eq:w}
 w=\sqrt{\frac{1-27z}{1+z}},\qquad z=\frac{1-w^2}{27+w^2}.
\end{equation}
The points $0,\rho$ correspond to $w=1,0$. Define
\begin{align}\label{eq:PTJ}
 T(w)&=15w^2+96w+85,\notag\\
 P(w)&=\left(\frac{7(w^2+27)}{T(w)}\right)^{1/4},\notag\\
 J(w)&=\frac{1728(1-w)^7(1+w)}{(w^2+27)T(w)^3},\qquad
 F(x)=\hyp{1/12}{5/12}{1}{x}.
\end{align}
Fourth roots are positive on $[0,1]$.
\begin{proposition}\label{prop:pullback}
As analytic germs near $w=1$, and throughout $0\le w\le1$,
\begin{equation}\label{eq:GF}
 A\!\left(\frac{1-w^2}{27+w^2}\right)=P(w)F(J(w)).
\end{equation}
\end{proposition}
\begin{proof}
The equation for $y(w)=A(z(w))$ is $y''+p y'+q y=0$, where
\begin{equation}\label{eq:wpq}
 p=\frac{2w}{(w-1)(w+1)},\qquad
 q=-\frac{16(w^2-15)}{(w-1)(w+1)(w^2+27)^2}.
\end{equation}
Let $h=P'/P=\tfrac14(2w/(w^2+27)-T'/T)$. The hypergeometric equation
\cite[\S15.10]{dlmf} is
\begin{equation}\label{eq:hypODE}
 x(1-x)F''+(1-\tfrac32x)F'-\tfrac5{144}F=0.
\end{equation}
Substitution of $y=PF(J)$, followed by eliminating $F''$, reduces the
claim to two rational identities:
\begin{align}\label{eq:certs}
 J''+(2h+p)J'+(J')^2\frac{3J/2-1}{J(1-J)}&=0,\notag\\
 h'+h^2+ph+q+\frac5{144}\frac{(J')^2}{J(1-J)}&=0.
\end{align}
Clearing denominators and expanding proves both; the attached code
checks them identically in $\Q(w)$. This is an exact certificate rather
than a fit to finitely many coefficients.
At $w=1$, $P=1$, $J=0$, and $z'(1)\ne0$. The right side is analytic
there and normalized to $1$. Uniqueness at $z=0$ proves the germ identity.
For $0\le w<1$,
\[
\frac{J'}J=-\frac7{1-w}+\frac1{1+w}-\frac{2w}{w^2+27}
                       -3\frac{30w+96}{T(w)}<0.
\]
Thus $0\le J(w)\le J(0)=64/85^3<1$, and every expression continues
regularly along the interval, proving the full claim.
\end{proof}
Differentiation at a singular $z$-value has now become ordinary
differentiation at the regular parameter value $w=0$.

\section{Evaluating the connection constant}
\subsection{An explicit hypergeometric expression}
Put
\begin{equation}\label{eq:q0}
 q_0=\frac{64}{85^3},\qquad P_0=\left(\frac{189}{85}\right)^{1/4}.
\end{equation}
As $s\to0$,
\[
 w=\sqrt{\frac{27s}{28-s}}
   =\sqrt{\frac{27}{28}}\sqrt s\,(1+O(s)).
\]
At $w=0$ we have $P'/P_0=-24/85$ and $J'/q_0=-798/85$. Hence
\begin{equation}\label{eq:B}
 B=-\sqrt{\frac{27}{28}}\frac{P_0}{85}
           \bigl(24F(q_0)+798q_0F'(q_0)\bigr)<0.
\end{equation}
Positivity follows from the positive hypergeometric coefficients.
Coefficient transfer for $\sqrt{1-27z}$ and
$\Gamma(-1/2)=-2\sqrt\pi$ gives
\begin{equation}\label{eq:Chyper}
 C=-\frac{B}{2\sqrt\pi}
 =\sqrt{\frac{27}{28}}\frac{P_0}{170\sqrt\pi}
      \bigl(24F(q_0)+798q_0F'(q_0)\bigr).
\end{equation}
This already gives an explicit positive constant. The following
special-value reductions explain its unexpectedly simple form.

\subsection{Ramanujan's identity removes the derivative}
The relevant Clausen identity is
\begin{equation}\label{eq:Clausen}
 F(q)^2=H(q):=\sum_{n\ge0}
      \frac{(1/6)_n(1/2)_n(5/6)_n}{(n!)^3}q^n.
\end{equation}
It follows directly from differential equations: the symmetric square
of \eqref{eq:hypODE} is
\[
 [\tht^3-q(\tht+1/6)(\tht+1/2)(\tht+5/6)]H=0,
 \qquad\tht=q\frac{d}{dq},
\]
whose analytic solution with $H(0)=1$ is unique. The rational
coefficients of this symmetric-square identity are checked in the code.

The classical level-seven Ramanujan identity, reproduced and obtained
by the construction in \cite[Eq.~(98)]{hemmecke}, is
\begin{equation}\label{eq:Ramanujan}
 \sum_{n\ge0}(133n+8)
   \frac{(1/6)_n(1/2)_n(5/6)_n}{(n!)^3}q_0^n
      =\frac{85\sqrt{255}}{54\pi}.
\end{equation}
The source uses the coefficient $(6n)!/((3n)!(n!)^3)$, which is $1728^n$
times the coefficient here, and $1728/255^3=64/85^3$. This normalization
check is essential. Equation \eqref{eq:Ramanujan} is a cited established
identity, not an equality inferred from numerical recognition.
Using \eqref{eq:Clausen}, it becomes
\[
 8F(q_0)^2+266q_0F(q_0)F'(q_0)=\frac{85\sqrt{255}}{54\pi}.
\]
Multiply by $3/F(q_0)$ and substitute into \eqref{eq:Chyper}. Since
$\Ast=P_0F(q_0)$ by Proposition~\ref{prop:pullback}, this yields
\begin{equation}\label{eq:Creciprocal}
 C=\frac{3\sqrt3}{8\pi^{3/2}P_0F(q_0)}
   =\frac{3\sqrt3}{8\pi^{3/2}\Ast}.
\end{equation}

\subsection{An elliptic period supplies the gamma product}
Let
\begin{equation}\label{eq:m7}
 m_7=\frac{8-3\sqrt7}{16},\qquad k_7=\sqrt{m_7},\qquad
 K(k)=\int_0^{\pi/2}\frac{d\theta}{\sqrt{1-k^2\sin^2\theta}}.
\end{equation}
Here $K$ takes the \emph{modulus} $k$, not the parameter $m=k^2$.
The classical singular-period evaluation is
\begin{equation}\label{eq:Kgamma}
 K(k_7)=\frac{\G}{4\pi7^{1/4}}.
\end{equation}
This is an established complex-multiplication special value
\cite{borweinzucker,beta}, not a new identity claimed here.
Appendix~\ref{app:beta} explicitly reduces the beta-function form in
Zucker's table to this normalization.

The required hypergeometric transformation is
\begin{equation}\label{eq:elliptictransform}
 F\!\left(\frac{27m^2(1-m)^2}{4(1-m+m^2)^3}\right)
 =(1-m+m^2)^{1/4}\hyp{1/2}{1/2}{1}{m}.
\end{equation}
For a direct verification put $Q(m)=1-m+m^2$ and
$j(m)=27m^2(1-m)^2/(4Q(m)^3)$. Then $Q^{-1/4}F(j)$ satisfies
\[
 m(1-m)y''+(1-2m)y'-\tfrac14y=0.
\]
Substituting and clearing denominators proves two rational identities,
exactly as in \eqref{eq:certs}. Both solutions are analytic at $0$ and
equal to $1$, proving the transformation near $0$. Continuation on
$[0,m_7]$ proves the value used here. The program verifies both identities.

At $m=m_7$,
\[
 m_7(1-m_7)=1/256,\qquad Q(m_7)=255/256,\qquad j(m_7)=q_0.
\]
Euler's integral gives
${}_2F_1(1/2,1/2;1;m_7)=2K(k_7)/\pi$. Therefore
\begin{align*}
 \Ast=P_0F(q_0)
 &=\left(\frac{189}{85}\frac{255}{256}\right)^{1/4}
                       \frac{2K(k_7)}\pi\\
 &=\frac{3\,7^{1/4}}4\frac2\pi
                       \frac{\G}{4\pi7^{1/4}}
 =\frac{3\G}{8\pi^2}.
\end{align*}
Combining this with \eqref{eq:Creciprocal} proves $C=\sqrt{3\pi}/\G$.
Since the coefficients are positive, the finite limit at $\rho$ also
gives the boundary sum \eqref{eq:boundarysum}, by monotone convergence.

\begin{remark}[Dependencies of the exact evaluation]
The two non-elementary special values are \eqref{eq:Ramanujan} and
\eqref{eq:Kgamma}; both are classical inputs. The recurrence
manipulations, rational pullbacks, connection-coefficient extraction,
and asymptotic deductions are proved here. We do not purport to give a
new derivation of Chowla--Selberg theory or of Ramanujan's general
$1/\pi$ formulas.
\end{remark}

\section{All orders: existence, remainders, and coefficients}
\subsection{The convergent local expansion}
Write $V(s)=\sum_{j\ge0}v_js^j$, $v_0=1$, $v_{-1}=0$.
Substitution of $s^{1/2}V(s)$ into \eqref{eq:sODE} gives
\begin{equation}\label{eq:vrec}
\begin{split}
28j(j+\tfrac12)v_j={}&
 \left(29(j-\tfrac12)^2+14(j-\tfrac12)+\frac{20}{9}\right)v_{j-1}\\
 &-\left((j-\tfrac32)^2+(j-\tfrac32)+\frac29\right)v_{j-2},\qquad j\ge1.
\end{split}
\end{equation}
Thus all $v_j$ are rational, with $v_1=593/1512$. Frobenius theory makes
this a \emph{convergent local} power series, not merely a formal solution
at infinity.
For nonintegral $\alpha$, the exact coefficient formula is
\begin{equation}\label{eq:transferexact}
 [z^n](1-27z)^\alpha
 =27^n\frac{\Gamma(n-\alpha)}{\Gamma(-\alpha)\Gamma(n+1)}.
\end{equation}
Apply it to $B\sum_{j=0}^M v_j(1-27z)^{j+1/2}$ and expand each gamma
ratio \cite[\S5.11]{dlmf}. Singularity analysis in the dented disk from
Section~3 transfers the remainder. The locally analytic part contributes
no algebraic terms from this singularity. Equivalently one may subtract
arbitrarily many Taylor terms of $U$ before applying the transfer theorem.
The first omitted nonintegral power is $(1-27z)^{M+3/2}$, and its
coefficient is $O(27^nn^{-M-5/2})$. This proves
\eqref{eq:dominant} for every fixed $M$.

The first correction is transparent: the gamma ratio for $\alpha=1/2$
contributes $3/(8n)$; the $s^{3/2}$ term contributes $-3v_1/(2n)$,
relative to the leading term. Hence
\[
 c_1=\frac38-\frac32\frac{593}{1512}=-\frac{215}{1008}.
\]
No numerical fitting of the coefficients is needed.

\subsection{A triangular recurrence at infinity}
For efficient computation insert $u_n=r^nn^{-3/2}S(1/n)$ in the
recurrence, with $r\in\{27,-1\}$. The corresponding formal operator is
\begin{equation}\label{eq:formalE}
\begin{split}
 \mathcal E_r[S](t)={}&r(1+t)^{1/2}S\!\left(\frac{t}{1+t}\right)
 -(26+13t+2t^2)S(t)\\
 &-\frac{27-27t+6t^2}{r}(1-t)^{-3/2}
                  S\!\left(\frac{t}{1-t}\right).
\end{split}
\end{equation}
The characteristic equation is $r^2-26r-27=0$, with roots $27,-1$.
\begin{proposition}[Exact coefficient algorithm]\label{prop:coeffalgo}
For either root there is a unique formal series
$S_r(t)=1+\sum_{j\ge1}\gamma_jt^j$ with $\mathcal E_r[S_r]=0$.
Writing $S_{m-1}=\sum_{j=0}^{m-1}\gamma_jt^j$, one has
\begin{equation}\label{eq:triangular}
 \gamma_m=\frac{[t^{m+1}]\mathcal E_r[S_{m-1}](t)}{m(r+27/r)},\qquad m\ge1.
\end{equation}
All coefficients are rational; the denominators in the recursion are
$28m$ and $-28m$, respectively.
\end{proposition}
\begin{proof}
A newly inserted $\gamma_mt^m$ contributes zero in order $m$, by the
characteristic equation. Its coefficient in order $m+1$ is
\[
 r(\tfrac12-m)-13-\frac{27}{r}(\tfrac12+m)=-m(r+27/r).
\]
This is nonzero, giving the assertion inductively. At the first two
orders the normalization is compatible with both characteristic roots.
\end{proof}
For $r=27$, the analytic existence proof above identifies this formal
series with the asymptotic series of $a_n/C$. For $r=-1$, Section~8 will
construct a genuine solution realizing the series. Formal coefficient
generation alone is not used to assert existence of an asymptotic
expansion.

\subsection{Logarithms, ratios, and convexity}
Let $\log S_{27}(t)=\sum_{j\ge1}\ell_jt^j$. Then
\begin{equation}\label{eq:logexp}
 \log a_n=n\log27-\frac32\log n+\log C
             +\sum_{j=1}^M\frac{\ell_j}{n^j}+O_M(n^{-M-1}),
\end{equation}
with
\begin{equation}\label{eq:ellvalues}
 \ell_1=-\frac{215}{1008},\quad\ell_2=-\frac{85}{3136},\quad
 \ell_3=\frac{99905}{85349376},\quad\ell_4=\frac{1657}{2458624}.
\end{equation}
In particular,
\begin{equation}\label{eq:ratio}
 \frac{a_{n+1}}{a_n}
 =27\left(1-\frac{3}{2n}+\frac{2105}{1008n^2}+O(n^{-3})\right).
\end{equation}
The sequence is strictly increasing, since the coefficient of $a_n$ in
\eqref{eq:rec} already exceeds $1$ for $n\ge1$, and $a_1>a_0$.
Taking sufficiently many terms before forming second differences gives
\[
 \log a_{n+1}-2\log a_n+\log a_{n-1}
 =\frac{3}{2n^2}+O(n^{-3})>0
\]
for sufficiently large $n$. Thus it is eventually strictly log-convex.
No explicit smallest threshold is asserted.

\section{A continued fraction and exact rational bracketing}
\subsection{The discrete Wronskian}
Let $b_n$ be as in Theorem~\ref{th:companion} and set
$W_n=a_nb_{n+1}-a_{n+1}b_n$. Eliminating the two forward terms gives
\[
 W_n=-\frac{3(3n-1)(3n-2)}{(n+1)^2}W_{n-1},\qquad W_0=1.
\]
Grouping factors by residue modulo $3$ yields the exact formula
\begin{equation}\label{eq:Wexact}
 W_n=(-1)^n\frac{(3n)!}{n!(n+1)!^2}.
\end{equation}
Consequently
\begin{equation}\label{eq:delta}
 \delta_n:=\frac{b_{n+1}}{a_{n+1}}-\frac{b_n}{a_n}
 =(-1)^n\frac{(3n)!}{n!(n+1)!^2a_na_{n+1}}.
\end{equation}
Its magnitude decreases strictly because
\[
 \left|\frac{\delta_{n+1}}{\delta_n}\right|
 =\frac{3(3n+2)(3n+1)a_n}{(n+2)^2a_{n+2}}<1.
\]
The numerator is one positive summand in the recurrence for the
denominator. The ratio tends to $1/27$ by Theorem~\ref{th:main}.
Thus $\sum\delta_n$ converges absolutely, and its alternating tails
prove the existence of $L$ and the brackets \eqref{eq:bracket}.
Even approximants increase and odd approximants decrease; consecutive
ones enclose $L$ with the exact width
\begin{equation}\label{eq:width}
 \left|\frac{b_{n+1}}{a_{n+1}}-\frac{b_n}{a_n}\right|
 =\frac{(3n)!}{n!(n+1)!^2a_na_{n+1}}.
\end{equation}
No unknown constants enter this stopping criterion.

\subsection{Polynomial continued-fraction coefficients}
The rescaled sequences $\mathcal H_n=(n!)^2a_n$ and
$\mathcal K_n=(n!)^2b_n$ satisfy
\begin{equation}\label{eq:integerrec}
 u_{n+1}=P_nu_n+Q_nu_{n-1},\qquad
 P_n=26n^2+13n+2,\quad Q_n=3n^2(3n-1)(3n-2),
\end{equation}
with initial pairs $(1,2)$ and $(0,1)$. Their integrality follows directly
from this recurrence, independently of integrality of the unscaled
sequence. The continuant recurrence gives
\begin{equation}\label{eq:CF}
 L=\cfrac1{2+\cfrac6{41+\cfrac{240}{132+\cfrac{1512}{275+\ddots}}}}.
\end{equation}
The convergent with $n$ denominator levels is
$\mathcal K_n/\mathcal H_n=b_n/a_n$.

\begin{remark}[The initial relation matters]
The companion satisfies the recurrence only for $n\ge1$;
$b_1=2b_0$ is false. Its generating function
$\mathcal B(z)=\sum b_nz^n$ therefore satisfies \eqref{eq:ODE} with right
side $1$, not $0$. The same holds for $\mathcal B-LA$.
These must not be called second analytic solutions of the homogeneous
equation at $0$. The two-dimensional solution space here consists of
sequences satisfying the recurrence for $n\ge1$.
\end{remark}

\section{The recessive scale and its full expansion}
Stirling's formula in \eqref{eq:Wexact} gives
\begin{equation}\label{eq:Wasym}
 W_n=(-1)^n\frac{\sqrt3}{2\pi}\frac{27^n}{n^3}(1+O(n^{-1})).
\end{equation}
Together with the dominant asymptotic this yields
\begin{equation}\label{eq:deltaasym}
 \delta_n=(-1)^n\frac{\sqrt3}{54\pi C^2}\,27^{-n}(1+O(n^{-1})).
\end{equation}
The exact tail identities are
\begin{equation}\label{eq:etail}
 L-\frac{b_n}{a_n}=\sum_{j\ge0}\delta_{n+j},\qquad
 e_n=-a_n\sum_{j\ge0}\delta_{n+j}.
\end{equation}
Since $\sum_{j\ge0}(-1/27)^j=27/28$, the leading terms of
\eqref{eq:CFerror} and \eqref{eq:minimal} follow. The recessive amplitude
can also be written
\begin{equation}\label{eq:Dgamma}
 D_{\min}=\frac{\sqrt3}{56\pi C}=\frac{\G}{56\pi^{3/2}}.
\end{equation}

\subsection{A remainder proof for the geometric tail}
The full Stirling expansion and Theorem~\ref{th:main} imply, for each
fixed $M$,
\[
 \delta_n=(-1)^nK_0\,27^{-n}
 \left(\sum_{k=0}^M\eta_kn^{-k}+O_M(n^{-M-1})\right),
 \qquad K_0=\frac{\sqrt3}{54\pi C^2},\quad\eta_0=1,
\]
where all $\eta_k$ are rational. The error bound is valid for all
sufficiently large integer arguments, hence also at $n+j$ uniformly in
$j\ge0$, with a bound by a constant times $27^{-j}n^{-M-1}$ after
factoring out $27^{-n}$.
Expand $(n+j)^{-k}$ for $0\le j\le n/2$ by Taylor's formula, retaining
its bounded remainder. The range $j>n/2$ is exponentially small, even
with any fixed polynomial weight. Every coefficient involves moments
\[
 \sum_{j\ge0}j^h(-1/27)^j,
\]
which are rational: apply $(q\,d/dq)^h$ to $1/(1-q)$ and set $q=-1/27$.
Thus the tail in \eqref{eq:etail} has an expansion to every fixed order
with a controlled error and rational normalized coefficients.
Multiplication by the expansion of $a_n$ proves existence of
\eqref{eq:minimal}. Since $e_n$ satisfies the recurrence,
Proposition~\ref{prop:coeffalgo} with $r=-1$ proves that its coefficients
are precisely the $d_j$ already given. This completes the proof of
Theorem~\ref{th:companion}.

\subsection{What the two-scale expansion does and does not say}
Because $W_n\ne0$, $a_n$ and $e_n$ form an exact basis of the recurrence
solutions for $n\ge1$. Every solution has a unique form
$u_n=\alpha a_n+\beta e_n$. In this specified basis its two asymptotic
scales are
\begin{equation}\label{eq:twoscale}
\begin{split}
 u_n\sim{}&\alpha C\,27^nn^{-3/2}\sum_{j\ge0}c_jn^{-j}\\
          &+\beta(-1)^{n+1}D_{\min}n^{-3/2}\sum_{j\ge0}d_jn^{-j}.
\end{split}
\end{equation}
Each exact basis sequence has its own fixed-order expansion. A finite
truncation of the dominant expansion generally has an error far larger
than the entire recessive term. Equation \eqref{eq:twoscale} is therefore
not a claim that two finite sums resolve exponentially small effects in
a generic solution. A canonical exponentially improved expansion of
$a_n$ alone would require additional summation or Stokes data; that
question is left open.

\section{Convolution powers and A183204}

For a fixed positive integer $m$, let
\[
 a_n^{\langle m\rangle}=[z^n]A(z)^m
 =\sum_{n_1+\cdots+n_m=n}a_{n_1}\cdots a_{n_m}.
\]
The exact boundary value controls this family as well as the original
sequence.

\begin{theorem}[Convolution asymptotics]\label{th:convolution}
For every fixed $m\ge1$,
\begin{equation}\label{eq:convmain}
 a_n^{\langle m\rangle}
 =m\Ast^{m-1}C\frac{27^n}{n^{3/2}}
 \left(1+\frac{c_1^{\langle m\rangle}}n+O_m(n^{-2})\right),
\end{equation}
where
\begin{equation}\label{eq:convc1}
 c_1^{\langle m\rangle}
 =-\frac{215}{1008}-\frac5{21}(m-1)
   -\pi\frac{C^2}{\Ast^2}(m-1)(m-2).
\end{equation}
There is a computable expansion to every fixed order, and all its
normalized coefficients lie in $\Q(\pi C^2/\Ast^2)$.
\end{theorem}

\begin{proof}
The analytic part of \eqref{eq:puiseux} satisfies
\[
 U(s)=\Ast\left(1+\frac{10}{63}s+O(s^2)\right),
\]
as the constant coefficient of \eqref{eq:sODE} shows. Its subsequent
normalized coefficients are rational. The singular part is
$B\sqrt{s}(1+v_1s+\cdots)$, with $B=-2\sqrt\pi C$ and
$v_1=593/1512$. In the $m$th power the coefficient of $\sqrt{s}$ is
$m\Ast^{m-1}B$. Relative to that coefficient, the coefficient of
$s^{3/2}$ is
\[
 v_1+(m-1)\frac{10}{63}
       +\frac{(m-1)(m-2)}6\frac{B^2}{\Ast^2}.
\]
Transfer as in Section~6 gives \eqref{eq:convc1}. At higher orders only
odd powers of $B$ multiply half-integral powers of $s$. After division
by the leading singular coefficient, the remaining expressions are
polynomials in $B^2/\Ast^2=4\pi C^2/\Ast^2$ with rational coefficients.
The convergent Puiseux expansion and the same dented domain provide the
fixed-order remainder estimates.
\end{proof}

For $m=2$ this is A183204, and the first two terms simplify to
\begin{equation}\label{eq:183204}
 [z^n]A(z)^2=
 \frac{3\sqrt3}{4\pi^{3/2}}\frac{27^n}{n^{3/2}}
 \left(1-\frac{65}{144n}+O(n^{-2})\right).
\end{equation}
The leading formula is already recorded in its OEIS entry
\cite{oeis183204}; it is recovered here, not claimed as new.
The cancellation of gamma factors at this convolution power is a useful
check on the evaluation of $C$.

For an algebraic check of the generating-function identification, the
symmetric-square equation of \eqref{eq:ODE} is
\begin{equation}\label{eq:squareODE}
\begin{split}
 (-z^2+26z^3+27z^4)Y'''&+(-3z+117z^2+162z^3)Y''\\
 &+(-1+86z+186z^2)Y'+(4+24z)Y=0.
\end{split}
\end{equation}
This is the equation recorded for A183204, and its normalized analytic
solution is unique. The program verifies this symmetric-square identity
and checks the convolution against
\[
 \sum_{k=0}^{n}\binom nk^2\binom{n+k}{n}\binom{2n-k}{n}
\]
for $0\le n\le40$. The finite check supplements the differential-equation
argument; it does not replace it.

\section{Inversion and the integer-rounding boundary}

Put $\lambda=\log27$ and $\beta=3/2$. The leading smooth model is
$Y=C e^{\lambda x}x^{-\beta}$. Its large positive solution is
\begin{equation}\label{eq:x0}
 x_0(Y)=-\frac{\beta}{\lambda}
 W_{-1}\!\left(-\frac{\lambda}{\beta}(C/Y)^{1/\beta}\right),
\end{equation}
where $W_{-1}$ is the real negative branch of Lambert's function.
The principal branch selects the small solution and is inappropriate
for inversion of the large-$n$ sequence \cite[\S4.13]{dlmf}.

\subsection{A triangular inverse expansion}

Seek
\begin{equation}\label{eq:inverse}
 x_M(Y)=x_0+\sum_{j=1}^{M}\frac{q_j}{x_0^j}.
\end{equation}
Writing $\delta=x-x_0$, the formal equation determining the coefficients
is
\begin{equation}\label{eq:inverseeq}
 \lambda\delta-\beta\log(1+\delta/x_0)
       +\sum_{k\ge1}\ell_k(x_0+\delta)^{-k}=0.
\end{equation}
The newly introduced coefficient $q_j$ has multiplier $\lambda$, so
this is triangular at every order. The first coefficients are
\begin{align}\label{eq:qvalues}
 q_1&=-\frac{\ell_1}{\lambda}
       =\frac{215}{1008\lambda},\notag\\
 q_2&=\frac{\beta q_1-\ell_2}{\lambda}
       =\frac{215}{672\lambda^2}+\frac{85}{3136\lambda},\notag\\
 q_3&=\frac{\beta q_2+\ell_1q_1-\ell_3}{\lambda}.
\end{align}
All $q_j$ belong to $\Q(\lambda)$.

An expansion involving only logarithms follows by putting
$X=\log(Y/C)$:
\begin{equation}\label{eq:pureloginverse}
\begin{split}
 x(Y)={}&\frac1\lambda\left\{
 X+\beta\log\frac X\lambda
 +\frac{\beta^2\log(X/\lambda)-\lambda\ell_1}{X}\right\}\\
 &\hspace{2em}+O\!\left(\frac{(\log X)^2}{X^2}\right).
\end{split}
\end{equation}
Here $x(Y)$ denotes the asymptotic inverse scale, not a previously
specified analytic interpolation of the integer sequence. Formula
\eqref{eq:x0} is often more accurate numerically because it retains the
entire leading model instead of truncating its logarithmic inversion.

\subsection{A rigorous discrete statement}

Define
\[
 N(Y)=\min\{n\ge0:a_n\ge Y\}.
\]
The sequence is strictly increasing and unbounded, so this is well-defined
for $Y>1$. An inverse approximation cannot uniformly remove the
ambiguity when its value is exceptionally close to an integer.

\begin{proposition}[Rounding brackets]\label{prop:rounding}
For every fixed $M\ge0$ there exist $K_M>0$ and $Y_M$ such that, for
$Y\ge Y_M$,
\begin{equation}\label{eq:rounding}
 \left\lceil x_M-K_Mx_M^{-M-1}\right\rceil
 \le N(Y)\le
 \left\lceil x_M+K_Mx_M^{-M-1}\right\rceil.
\end{equation}
When the two ceilings coincide, they determine the discrete inverse
exactly. For $M=0$, $x_M$ means $x_0$.
\end{proposition}

\begin{proof}
Use the smooth finite model
\[
 g_M(x)=\log C+\lambda x-\beta\log x+
                      \sum_{k=1}^{M}\ell_kx^{-k}.
\]
The all-orders theorem gives
$\log a_n=g_M(n)+O_M(n^{-M-1})$. Also
$g_M'(x)=\lambda+O(x^{-1})$, so it is bounded away from zero for large
$x$. Let $t_M$ be the large solution of $g_M(t_M)=\log Y$.
Formal substitution of \eqref{eq:inverse} into $g_M$, with ordinary
finite Taylor remainders, gives
$g_M(x_M)-\log Y=O_M(x_M^{-M-1})$. The mean value theorem therefore
implies $t_M-x_M=O_M(x_M^{-M-1})$.

For integers near $t_M$, comparison of $g_M(n)$ with $g_M(t_M)$,
using its derivative lower bound and the error in $\log a_n$, shows
that integers sufficiently below $t_M$ cannot reach $Y$, whereas
integers sufficiently above it do reach $Y$. Enlarge the constant to
absorb $t_M-x_M$, and use monotonicity of $a_n$ for all remaining
integers. This gives the two ceilings in \eqref{eq:rounding}.
\end{proof}

The proposition asserts the existence of fixed-order constants, not
explicit numerical values of $K_M$ or $Y_M$. Computing certified,
practical values is one of the research directions below. No claim
that $N(Y)=\lceil x_M(Y)\rceil$ holds uniformly is made.

\section{Reproducible computations and certified constants}

Build and run instructions are in the accompanying \texttt{README.md};
the source audit is in \texttt{SOURCE\_AUDIT.md}. The PDF can be built
without running Python, since its numerical tables are part of the TeX
source. The accompanying program uses exact integers and rational functions for
its algebraic checks. Arbitrary-precision floating-point arithmetic is
used separately for numerical diagnostics. It requires only Python,
SymPy, and mpmath, and makes no network requests.

\subsection{Certificates and reproducible output}

The following assertions were executed successfully: eight rational
identities for the two pullbacks and two symmetric squares; coefficient
recurrences through order $12$ for both characteristic roots; exact
recurrence evaluation through $n=1000$; $100$ discrete Wronskian
identities; and $41$ comparisons with the binomial sum for A183204.
The division in the $a_n$ recurrence is exact at every tested index.
This last computation is a finite integrality check, not an independent
proof of integrality for all indices.

The constants are
\begin{align*}
 C&=0.09552230526812671465130791078702962567279466665100\ldots,\\
 \Ast&=1.22113109746756666928422458568599346237582705637935\ldots,\\
 L&=0.46721176888437312394211005141309126172410854995971\ldots,\\
 D_{\min}&=0.10306659600794607491874204424062776206902297866964\ldots.
\end{align*}
The first displayed value is not inferred by integer-relation searching:
its gamma formula is proved in Section~5, and the following interval
method supplies independent certified decimal evaluation.

\subsection{A rational enclosure for the asymptotic constant}

Write $F(q_0)=\sum_{j\ge0}h_j$. Its positive terms satisfy
\[
 \frac{h_{j+1}}{h_j}
 =q_0\frac{(12j+1)(12j+5)}{144(j+1)^2}<q_0<1.
\]
Hence, for $S_N=\sum_{j=0}^{N-1}h_j$,
\begin{equation}\label{eq:Fbound}
 S_N\le F(q_0)\le S_N+\frac{h_N}{1-q_0}.
\end{equation}
Every quantity in this enclosure is rational. The implementation bounds
$\pi$ using Machin's identity
\[
 \pi=16\arctan(1/5)-4\arctan(1/239)
\]
and exact alternating-series bounds. Square and fourth roots are
bracketed by rational numbers using integer root extraction on a fixed
decimal grid. Combining positive upper and lower bounds in
\eqref{eq:Creciprocal} gives a rational enclosure for $C$ of width less
than $10^{-120}$. Outward rounding produces the following interval of
width $10^{-100}$; each row continues on its second line.

\begin{center}\small\ttfamily
\begin{tabular}{ll}
Lower:&0.09552230526812671465130791078702962567279466665100\\
      &71798669948234917659264524531457675812093972181446\\[3pt]
Upper:&0.09552230526812671465130791078702962567279466665100\\
      &71798669948234917659264524531457675812093972181447
\end{tabular}
\end{center}

For $L$ the enclosure uses no transcendental functions at all:
$b_{90}/a_{90}<L<b_{91}/a_{91}$, with exact rational endpoints whose
difference is less than $10^{-125}$. The full endpoints and their
100-place decimal enclosures are stored in
\texttt{data/verification.json}. The JSON also records every coefficient
through order $12$ and the zero values of all eight rational certificates.

\subsection{Numerical diagnostics}

For Table~\ref{tab:dominant}, the relative error is
\[
 \frac{C27^nn^{-3/2}\sum_{j=0}^{M}c_jn^{-j}}{a_n}-1.
\]
The values are rounded floating-point diagnostics, not new rigorous
upper bounds on the asymptotic remainder. All were computed with
190 decimal digits of working precision.

\begin{table}[htbp]\centering\small
\caption{Dominant expansion: signed relative errors.}\label{tab:dominant}
\begin{tabular}{rllll}\toprule
$n$ & $M=0$ & $M=3$ & $M=6$ & $M=12$\\\midrule
20 & $1.079\times10^{-2}$ & $-1.143\times10^{-9}$ & $-5.413\times10^{-12}$ & $2.818\times10^{-19}$\\
100 & $2.138\times10^{-3}$ & $-2.449\times10^{-12}$ & $-6.609\times10^{-17}$ & $6.083\times10^{-28}$\\
1000 & $2.133\times10^{-4}$ & $-2.597\times10^{-16}$ & $-6.536\times10^{-24}$ & $6.727\times10^{-41}$\\
\bottomrule\end{tabular}\end{table}

The normalized recessive error in Table~\ref{tab:minimal} is
$\sum_{j=0}^{M}d_jn^{-j}-e_n/[(-1)^{n+1}D_{\min}n^{-3/2}]$.
The value of $L$ used in this diagnostic is the midpoint of the exact
$r_{150},r_{151}$ enclosure, where $r_n=b_n/a_n$. Its error remains
negligible after multiplication by $a_n$ at the displayed indices.

\begin{table}[htbp]\centering\small
\caption{Recessive expansion: signed normalized errors.}\label{tab:minimal}
\begin{tabular}{rlll}\toprule
$n$ & $M=0$ & $M=3$ & $M=6$\\\midrule
10 & $4.910\times10^{-2}$ & $1.721\times10^{-6}$ & $-1.532\times10^{-9}$\\
30 & $1.676\times10^{-2}$ & $2.361\times10^{-8}$ & $-6.532\times10^{-13}$\\
75 & $6.753\times10^{-3}$ & $6.240\times10^{-10}$ & $-1.036\times10^{-15}$\\
\bottomrule\end{tabular}\end{table}

Table~\ref{tab:inverse} tests inversion at $Y=a_n$. It reports the signed
continuous error $x_M(a_n)-n$, not an integer rounding error. The
uncorrected Lambert-$W$ scale is already accurate; the algebraic
corrections produce the predicted further improvement.

\begin{table}[htbp]\centering\small
\caption{Inverse expansion at exact sequence values.}\label{tab:inverse}
\begin{tabular}{rlll}\toprule
$n$ & $x_0(a_n)-n$ & $x_1(a_n)-n$ & $x_3(a_n)-n$\\\midrule
20 & $-3.332\times10^{-3}$ & $-9.581\times10^{-5}$ & $-6.315\times10^{-9}$\\
100 & $-6.509\times10^{-4}$ & $-3.780\times10^{-6}$ & $-1.064\times10^{-11}$\\
1000 & $-6.475\times10^{-5}$ & $-3.769\times10^{-8}$ & $-1.075\times10^{-15}$\\
\bottomrule\end{tabular}\end{table}

\section{Further research questions and topics}

The questions below are concrete extensions of the proved results.
They are not assertions that every item is globally open: before a
priority claim, the relevant modular-form and continued-fraction
literature should be checked in greater depth.

\subsection{Explicit remainder constants}
Theorems~\ref{th:main} and \ref{th:companion} give remainders at every
fixed order but not practical values of the constants. Construct
computable $K_M,N_M$ such that
\[
 \left|\frac{a_n n^{3/2}}{C27^n}-\sum_{j=0}^{M}c_jn^{-j}\right|
 \le K_M n^{-M-1}\qquad(n\ge N_M).
\]
One route is to bound the analytic Frobenius functions on a specified
complex neighborhood and make the transfer contour explicit. A second
is to construct upper and lower recurrence solutions from truncated
series plus a controlled correction. The latter may give smaller
thresholds and could turn the inverse brackets into a practical
certified algorithm.

\subsection{Late-order growth and optimal truncation}
Determine the large-$j$ behavior of $c_j$ and $d_j$. The proof of
fixed-order asymptotics does not imply convergence of the series in
$1/n$. A useful target is a sharp factorial-over-power law, with its
constant and phase, or a proof of a specified Gevrey bound. This would
identify an optimal truncation index and quantify what accuracy is
available before the recessive scale becomes relevant.

\subsection{Canonical exponentially improved asymptotics}
The exact basis $a_n,e_n$ removes ambiguity from
\eqref{eq:twoscale}, but it does not specify a canonical summation of the
dominant formal series. Develop a Borel or contour normalization and
compute the associated connection data between the singularities
$\rho$ and $-1$. The desired result is an error theorem resolving the
alternating contribution after optimal truncation, not merely a formal
addition of a second exponential scale.

\subsection{Recognizing the continued-fraction limit}
Find an exact special-value expression for $L$. The inhomogeneous
operator in the remark following \eqref{eq:CF} provides a starting point:
variation of parameters expresses its solutions through integrals
involving $A$ and algebraic functions. The condition that the
$\sqrt{1-27z}$ coefficient vanish selects $\mathcal B-LA$ and hence $L$.
A successful reduction might involve an elliptic integral, an iterated
modular integral, or a special $L$-value. No such identification is
proved in this report.

\subsection{Arithmetic strength of the approximants}
Does the continued fraction prove irrationality of $L$, or a useful
irrationality measure? The exponential error $27^{-n}$ alone is not
enough: the natural continuant denominators contain factorial growth.
The first step is to analyze the reduced denominators and common
factors of $\mathcal H_n$ and $\mathcal K_n$. Any arithmetic conclusion
must compare the error with the actual denominator after cancellation.
The general results in \cite{cohenzudilin} are relevant context, but
do not automatically apply to the specific limit constructed here.

\subsection{The arithmetic conjecture in the entry}
The OEIS entry also displays Lucas congruences for primes in the
specified residue classes modulo $7$ \cite{oeis279619}. The present
asymptotic argument neither proves nor disproves them. A separate
project should first audit subsequent work on the Chan--Cooper--Sica
conjecture and then exploit the modular parametrization or the
functional equation added to OEIS in 2026. This is a different target
from the asymptotic conjecture resolved here; analytic evidence should
not be substituted for an arithmetic proof.

\subsection{Coefficient fields for convolution powers}
Theorem~\ref{th:convolution} bounds the coefficient field by
$\Q(\pi C^2/\Ast^2)$. Determine exactly which convolution powers and
which correction orders eliminate the period ratio. The square is
special at its first correction, and its third-order differential
equation may expose cancellations more systematically. A symbolic
classification would clarify when a convolution inherits purely
rational normalized corrections and when new period combinations
necessarily appear.

\subsection{Noninteger powers and growing convolution order}
Extend the convolution theorem to a specified branch of $A(z)^u$ for
noninteger real or complex $u$. Zero avoidance and possible additional
singularities must be addressed rather than assumed. A separate
uniform problem lets the integer $m=m(n)$ grow. The fixed-$m$ proof
is then insufficient because powers of the analytic boundary part
and the singular part can compete. Determine the applicable uniform
ranges and the transition scales from the actual local expansion.

\subsection{Positive integral representations}
Find a direct integral or moment representation for $a_n$ that makes
positivity and possible stronger inequalities transparent. A compact
real support reflecting both characteristic roots would be especially
informative, but such a representation and positivity of its density
are not established here. Potential consequences include explicit
log-convexity thresholds, Hankel determinant signs, and alternative
error bounds for coefficient asymptotics.

\subsection{Automated connection-constant recognition}
For families of quadratic-coefficient recurrences with two distinct
characteristic roots, combine three exact procedures: derive the
Fuchsian equation, certify a rational hypergeometric pullback, and test
whether a singular endpoint is a known complex-multiplication value.
The intended output should include a symbolic certificate and a branch
check, not just a high-precision guess. A279619 is a compact test case
because the rational parameter turns the dominant endpoint into $w=0$.

\subsection{A direct modular proof of the product identity}
Can the product $C\Ast=3\sqrt3/(8\pi^{3/2})$ be obtained directly from a
single level-seven modular relation? The proof here combines a
Ramanujan identity with a period evaluation. A direct connection or
Wronskian argument on the modular side may unify the two normalizations
and generalize to other Ap\'ery-like equations. Such a proof could also
explain which sequences exhibit reciprocal boundary and singular
periods.

\subsection{Machine-checked certificates and analytic infrastructure}
A staged Lean development can begin with the recurrence, positivity,
integer continuants, and the exact Wronskian. These statements need only
algebra and induction. The rational identities in
\eqref{eq:certs} and the triangular coefficient algorithm form a second
self-contained layer. A later analytic layer must connect convergent
local solutions to coefficient asymptotics, including branch choices
and remainder estimates. The complex-multiplication identities should
be stated as explicit dependencies until they too are formalized.
This separation prevents an exact symbolic certificate from being
mistaken for a proof of analytic continuation or of a special value.

\section{Integration roadmap and conclusion}

A natural repository layout separates four components: sequence
recurrences and exact certificates; local analytic representation;
asymptotic coefficient calculus; and inverse expansions. The first and
third components can be tested independently of modular-form machinery.
The second supplies the analytic existence and connection constants
that formal manipulations alone cannot establish. The fourth should
retain an explicit distinction between a continuous inverse scale and
an integer threshold function.

No existing Lean declaration is assumed to certify the special-function
steps of this report. A useful first integration milestone is a
machine-checked discrete Wronskian with rational bracketing of $L$.
The next is a generic formal-series lemma giving
\eqref{eq:triangular}; both are substantially smaller than a full
formalization of the gamma constant. The supplied Python verification
program is a reproducible algebraic audit and a source of test data,
not a replacement for that formalization.

The principal outcome is a precise resolution of the asymptotic
statement displayed for A279619:
\[
 \boxed{\displaystyle
 a_n\sim\frac{\sqrt{3\pi}}{\Gamma(1/7)\Gamma(2/7)\Gamma(4/7)}
                    \frac{27^n}{n^{3/2}}.}
\]
The same argument supplies the boundary sum, every algebraic correction
with a fixed-order error bound, and a useful convolution calculus.
The independent companion construction reveals an alternating
recessive solution with an exact gamma-product amplitude. Inversion
then yields a Lambert-$W$ scale and rigorous rounding brackets. The
remaining questions concern sharper effective bounds, canonical
exponentially improved expansions, arithmetic properties, and broader
families, rather than the existence or value of the leading asymptotic
constant proved here.

\appendix
\section{Checking the classical period normalization}\label{app:beta}

The table in Zucker's singular-value notes \cite[Table~1]{beta} uses
$K[N]=K(k_N)$ with $K(k_N')/K(k_N)=\sqrt N$. At $N=7$ it gives
\begin{equation}\label{eq:betatable}
 K[7]=\frac{2^{5/7}7^{3/4}}{14}
 \frac{\beta(1/7)\beta(2/7)}{\beta(1/14)},\qquad
 k_7^2=\frac{8-3\sqrt7}{16},
\end{equation}
where $\beta(t)=\Gamma(t)^2/\Gamma(2t)$. We spell out the simplification
because confusing the elliptic modulus with its square would give an
incorrect constant.

Set $g_1=\Gamma(1/7)$, $g_2=\Gamma(2/7)$, $g_4=\Gamma(4/7)$, and
$g_{1/2}=\Gamma(1/14)$. Then
\[
 \frac{\beta(1/7)\beta(2/7)}{\beta(1/14)}
 =\frac{g_1^3g_2}{g_4g_{1/2}^2}.
\]
The gamma duplication formula at $1/14$ is
\[
 g_{1/2}g_4=2^{6/7}\sqrt\pi\,g_1.
\]
It follows that the beta quotient is
$g_1g_2g_4/(2^{12/7}\pi)=\G/(2^{12/7}\pi)$. Substitution into
\eqref{eq:betatable} gives
\[
 K[7]=\frac{7^{3/4}\G}{28\pi}
      =\frac{\G}{4\pi7^{1/4}},
\]
which is exactly \eqref{eq:Kgamma}. The classical table evaluation is
an external input; the algebraic conversion here checks that its
normalization matches the one used in the proof.

\section{A possible OEIS formula update}

The following is suggested mathematical content for an update after
independent review. No edit to OEIS or to the user's repository has
been made as part of this task.

\begin{resultbox}
Let $G=\Gamma(1/7)\Gamma(2/7)\Gamma(4/7)$. Then
\[
 a(n)\sim\frac{\sqrt{3\pi}}{G}\frac{27^n}{n^{3/2}},\qquad
 \sum_{n\ge0}\frac{a(n)}{27^n}=\frac{3G}{8\pi^2}.
\]
More precisely, for every fixed truncation order there is an asymptotic
expansion with a remainder of the next order. Its first corrections are
\[
 1-\frac{215}{1008n}-\frac{1265}{290304n^2}
   +\frac{4683055}{877879296n^3}+O(n^{-4}).
\]
The first three corrections were already calculated formally by
O'Brien \cite{obrien}.
\end{resultbox}

The separate Lucas-congruence statement is not affected by this report.

\clearpage
\begin{thebibliography}{99}
\raggedright

\bibitem{oeis279619}
The OEIS Foundation Inc., \emph{A279619: Expansion of the generating
function of A002652 in powers of the generating function of A279618}.
Entry by Lynette O'Brien, with later contributions; consulted October 1,
2026. The displayed entry revision was August 22, 2026.
\url{https://oeis.org/A279619/internal}.

\bibitem{oeis183204}
The OEIS Foundation Inc., \emph{A183204}, recurrence, differential
equation, generating-function relations, and asymptotic formula.
Consulted October 1, 2026.
\url{https://oeis.org/A183204/internal}.

\bibitem{obrien}
L.~O'Brien, \emph{Modular forms and two new integer sequences at level 7},
Master of Science thesis, Massey University, 2016.
Especially Theorem~6.1 and Section~8.4, pp.~63--67.
\url{https://oeis.org/A279619/a279619.pdf}.

\bibitem{hemmecke}
R.~Hemmecke, P.~Paule, and C.-S.~Radu,
\emph{Computer-assisted construction of Ramanujan--Sato series for
1 over $\pi$}, manuscript dated January 31, 2025.
Equation~(98), p.~32, is the level-seven identity used here.
\url{https://www3.risc.jku.at/publications/download/risc_7134/.RamanujanSatoOneOverPi.pdf}.

\bibitem{borweinzucker}
J.~M.~Borwein and I.~J.~Zucker,
``Fast evaluation of the gamma function for small rational fractions
using complete elliptic integrals of the first kind,''
\emph{IMA Journal of Numerical Analysis} \textbf{12} (1992), 519--526.
\url{https://academic.oup.com/imajna/article-abstract/12/4/519/690281}.

\bibitem{beta}
I.~J.~Zucker, \emph{Singular value theory}, exercise notes in
Jonathan Borwein's archived materials, Table~1, entry $K[7]$.
The notes refer to Borwein and Zucker's elliptic-integral evaluations;
consulted October 1, 2026.
\url{https://carmamaths.org/jon/Preprints/Books/EMA/Exercises/For%20others/K-beta.pdf}.

\bibitem{flajolet}
P.~Flajolet and R.~Sedgewick, \emph{Analytic Combinatorics}, Cambridge
University Press, 2009. Chapter~VI: singularity analysis and transfer
of local expansions.

\bibitem{dlmf}
F.~W.~J.~Olver et al., eds., \emph{NIST Digital Library of Mathematical
Functions}, Sections~4.13, 5.11, and 15.10: Lambert's $W$ function,
gamma asymptotics, and the hypergeometric differential equation.
Consulted October 1, 2026.
\url{https://dlmf.nist.gov/}.

\bibitem{cohenzudilin}
H.~Cohen and W.~Zudilin,
\emph{Continued fractions and irrationality measures for
Chowla--Selberg gamma quotients}, arXiv:2510.00215, 2025, revised
versions through 2026. Related work, not a theorem dependency here.
\url{https://arxiv.org/abs/2510.00215}.

\bibitem{proveit}
V.~Reshetnikov, \emph{ProveIt},
\texttt{Analysis/Transseries/README.md}, repository snapshot
\texttt{a5ad4ac2e2d4e642a6c015836a70492d40fcbc71}.
\url{https://github.com/VladimirReshetnikov/ProveIt/tree/main/Analysis/Transseries}.

\end{thebibliography}
\end{document}
